\documentclass[10pt,conference]{IEEEtran}
\usepackage[T1]{fontenc}
\usepackage[utf8]{inputenc}
\usepackage{amsmath,amssymb}
\usepackage{url}
\usepackage{booktabs}
\usepackage{hyperref}
\hypersetup{colorlinks=true,linkcolor=blue,urlcolor=blue}
\usepackage{enumitem}

\begin{document}

\title{Technical Disclosure: PURIQ Master Formula, the 13-Axis Yuyay Gate, and a
Khipu DAG Architecture for Provable Agentic AI}

\author{
\IEEEauthorblockN{Yachay (Evaluation \& Defense), on behalf of SZL Holdings}
\IEEEauthorblockA{SZL Holdings --- PURIQ Agentic Anatomy Layer\\
Defensive Publication. Disclosure date: 2026-06-01 (EDT).\\
Concept DOI chain: \href{https://doi.org/10.5281/zenodo.19944926}{10.5281/zenodo.19944926}}
}

\maketitle

\begin{abstract}
This document is a \emph{defensive technical disclosure} establishing a public,
timestamped record of the PURIQ agentic-AI architecture so that the methods herein
enter the prior art and cannot be exclusively claimed by third parties. We disclose:
(i) the PURIQ \emph{master action-selection formula} $P(x,t)$ that gates every agentic
action on a bounded utility aggregator, a 13-axis wisdom score, an exponential
safety-tripwire penalty, and a product of cryptographic receipt verifications;
(ii) the \emph{13-axis Yuyay gate}, a conjunctive admissibility test; and (iii) the
\emph{Khipu DAG}, a SHA-256 hash-chained receipt structure giving per-action provenance
and an $O(1)$ integrity tripwire. We anchor the disclosure to concrete, independently
checkable artifacts: a formal-methods corpus of 749 Lean declarations with 14 unique
axioms and 13 sorry-free (PROVED) theorems, 23 extracted mathematical formulas, a
deterministic replay-hash anchor (\texttt{bacf5443\ldots631fc5}), Zenodo DOIs, and a
load test demonstrating 5.35\,M receipt writes/min with a 100\% chain-integrity
tripwire. We state honestly which claims are PROVED, which remain open obligations
(\texttt{sorry}-tagged), and which provenance levels (SLSA L1) are asserted.
\end{abstract}

\begin{IEEEkeywords}
agentic AI, provenance, hash chain, formal verification, Lean, defensive publication,
action selection, AI governance, tool-use safety.
\end{IEEEkeywords}

\section{Field and Purpose}
The disclosed subject matter relates to governed, auditable agentic AI systems that
select and execute actions (tool calls, code edits, retrievals) under verifiable safety
and provenance constraints. The purpose of this publication is \emph{defensive}: to place
the architecture, formula, and data structures into the public prior art as of the
disclosure date, thereby preventing later third-party patents from foreclosing SZL's
freedom to operate, while SZL separately pursues a narrow set of patent claims
(Section~\ref{sec:patent}).

\section{The PURIQ Master Formula}
An agentic system maintains a context $x$ and, at time $t$, selects an action $a$ from a
bounded action space $\mathcal{A}$. PURIQ selects:
\begin{equation}
P(x,t) = \operatorname*{arg\,max}_{a \in \mathcal{A}}
\Big[\, \Lambda(x)\cdot \mathrm{Yuyay}_{13}(a)\cdot
e^{-\beta\, \mathrm{HUKLLA}(a)}\cdot \prod_i \mathrm{Khipu}_i(a) \,\Big]
\label{eq:master}
\end{equation}
where:
\begin{itemize}[leftmargin=1.2em]
\item $\Lambda(x)$ --- a \emph{Lambda Spine aggregator}: positive-homogeneous, bounded,
monotone. (Uniqueness of the stationary aggregator is asserted as \emph{Conjecture~1},
not a theorem --- stated honestly.)
\item $\mathrm{Yuyay}_{13}(a) \in [0,1]$ --- the 13-axis wisdom score
(Section~\ref{sec:yuyay}); \emph{conjunctive}: it is $0$ unless all 13 axes clear their
floors.
\item $\mathrm{HUKLLA}(a)$ --- count of tripwire violations $T01\ldots T10$, penalized
exponentially with coefficient $\beta$. As $\beta\!\to\!\infty$, any violating action's
score $\to 0$ (\emph{halting safety}).
\item $\prod_i \mathrm{Khipu}_i(a)$ --- product of receipt verifications
(Section~\ref{sec:khipu}); each factor is $1$ iff \texttt{chain\_verified} for that
receipt, else $0$. A single broken receipt zeroes the entire score
(\emph{chain-integrity required for non-zero score}).
\item $\mathcal{A}$ --- bounded action space (information-theoretically bounded by the
context, a Bekenstein-style cap on $|\mathcal{A}|$).
\end{itemize}

\paragraph{Disclosed invariants (proof obligations).}
(INV-1) $\Lambda$-boundedness and monotonicity; (INV-2) Yuyay conjunctive gating;
(INV-3) HUKLLA halting safety as $\beta$ grows; (INV-4) Khipu chain-integrity
necessity. We disclose these as the four invariants; in the SZL Lean corpus they are
stated as theorems, several PROVED and some still \texttt{sorry}-tagged. We do
\emph{not} claim all four are machine-proven; we claim the formula and the invariant
\emph{statements} as prior art.

\section{The 13-Axis Yuyay Gate}\label{sec:yuyay}
$\mathrm{Yuyay}_{13}$ partitions admissibility into 13 named axes evaluated as a
conjunction (logical AND) with per-axis floors:
\begin{itemize}[leftmargin=1.2em]
\item \textbf{2 sacred axes}, floor $\geq 0.95$ each;
\item \textbf{7 structural axes}, floor $\geq 0.90$ each;
\item \textbf{4 introspection axes}, cross-linked to HUKLLA tripwires
$T03, T04, T09, T10$.
\end{itemize}
The gate returns a score in $[0,1]$ that collapses to $0$ if \emph{any} axis fails its
floor, making the gate a hard admissibility filter rather than a soft average. The axis
configuration is anchored to the deterministic replay-hash
\texttt{bacf54434f1a3bf2d758b27a62d5fd580ca4c8d3b180693573eeebcaea631fc5}, so that any
re-evaluation of the same action under the same axis weights reproduces the same gate
decision bit-for-bit. We disclose the partition (2/7/4), the floors (0.95/0.90), and the
introspection-to-tripwire cross-link as prior art.

\section{The Khipu DAG: Hash-Chained Receipts}\label{sec:khipu}
Every action emits a \emph{Khipu receipt}. Receipts form a directed acyclic graph of
per-domain chains. For a chain head $h_{\text{prev}}$ and a canonical action payload
$m$, the next receipt hash is
\begin{equation}
h = \mathrm{SHA256}\big(h_{\text{prev}} \,\Vert\, \mathrm{canonical}(m)\big),
\label{eq:chain}
\end{equation}
with genesis $h_{\text{prev}} = \texttt{GENESIS}$. Verification recomputes
Eq.~\eqref{eq:chain} for every receipt and checks linkage; any mismatch (tampered hash
or broken \texttt{prev} pointer) is a tripwire $T01$ event, and by
Eq.~\eqref{eq:master} forces $\prod_i \mathrm{Khipu}_i(a)=0$. The DAG fans out across
independent organ chains (e.g.\ six SZL Spaces: a11oy, amaru, sentra, killinchu, rosie,
vessels), enabling parallel write with global verifiability.

\paragraph{Empirical anchor.} An in-process measurement of this construction wrote
$500{,}000$ receipts across six chains at a sustained
\textbf{5{,}349{,}190 receipts/min} (89{,}153/s), verified at 232{,}830/s, with
per-write latency p50/p95/p99 of $4.83/8.31/17.14\,\mu s$. Of $504$ deliberately
tampered receipts, the $T01$ tripwire caught \textbf{504/504 (100\%)}. These numbers are
disclosed as evidence of practical reduction to practice of the integrity mechanism.

\section{Anchoring Artifacts (timestamped prior art)}
The following independently verifiable artifacts anchor the disclosure date and content:
\begin{table}[h]
\centering
\caption{Disclosed anchoring artifacts}
\begin{tabular}{@{}ll@{}}
\toprule
Artifact & Value / Reference \\
\midrule
Concept DOI (thesis chain) & 10.5281/zenodo.19944926 \\
a11oy deposit DOI & 10.5281/zenodo.20451991 \\
lutar-lean (canonical cite) & 10.5281/zenodo.20434308 \\
Lean declarations & 749 \\
Unique axioms & 14 \\
PROVED (sorry-free) theorems & 13 \\
\texttt{sorry} obligations (raw) & 163 (honest; not hidden) \\
Extracted formulas & 23 \\
Replay-hash anchor & \texttt{bacf5443\ldots631fc5} \\
Lean corpus commit & \texttt{c7c0ba1} (lutar-lean) \\
Provenance level & SLSA L1 (honest; ``L3'' is not claimed) \\
\bottomrule
\end{tabular}
\end{table}

We state honestly: $\Lambda$-uniqueness is \emph{Conjecture~1}, not a theorem; the
signature factor is a hash-chain verifier (a DSSE signature binding is a placeholder, not
yet implemented); and 163 \texttt{sorry} obligations remain open in the corpus. These
honest caveats are themselves part of the disclosed record so that the prior art is
accurate and not overclaimed.

\section{Distinguishing Features Disclosed as Prior Art}
\begin{enumerate}[leftmargin=1.4em]
\item A single action-selection objective (Eq.~\ref{eq:master}) that
\emph{multiplicatively} composes a bounded utility aggregator, a conjunctive multi-axis
admissibility gate, an exponential tripwire penalty, and a product of cryptographic
receipt checks --- such that violation of \emph{any} factor zeroes the action's score.
\item A 13-axis (2/7/4) conjunctive gate with explicit floors and tripwire cross-links,
anchored to a deterministic replay hash for bit-reproducible decisions.
\item A SHA-256 DAG of per-organ receipt chains in which chain-integrity is a
\emph{necessary condition} for any action to be selectable, with an $O(1)$ per-receipt
verification tripwire.
\item Routing of inference by \emph{sovereignty/governance tier} and \emph{license
class} as a precondition embedded in the same gated objective.
\end{enumerate}

\section{Patent Strategy: File vs.\ Publish-Defensively}\label{sec:patent}
SZL's strategy separates a small number of \emph{narrow, novel} claims to \emph{file}
from a broad \emph{fence} to \emph{publish defensively} (this document):
\begin{itemize}[leftmargin=1.2em]
\item \textbf{FILE (3 claims).}
\textbf{(P-A) Sovereignty-Selectable Inference} --- selecting the inference provider by a
sovereignty/governance tier and license boundary as an enforced gate (strongest concrete
technical effect; regulatory wedge). \textbf{(P-B) Gate-Minted Capability Token} --- a
capability/authorization token minted only upon passing the 13-axis gate, bounding
downstream tool authority. \textbf{(P-C) Theorem-Bound Tool Output} --- binding a tool
output to a \emph{named, machine-proved} theorem from the 13-PROVED set, scoped narrowly
to avoid prior art on draft-model selection.
\item \textbf{PUBLISH DEFENSIVELY (this disclosure).} The broad master formula
(Eq.~\ref{eq:master}), the general ``route-by-sovereignty'' concept fence around P-A,
the Khipu DAG integrity mechanism, and the cross-provider speculative-decoding slice ---
which is folded into P-A's license-boundary claim rather than filed standalone, because
existing patents (e.g.\ US12229192B2, US20250384043A1) densely cover draft-model
selection and self-speculation, making a standalone filing high freedom-to-operate risk.
\end{itemize}
Sequencing: file P-A first; publish the broad fence simultaneously; hold P-B/P-C until a
signature/attestation (e.g.\ Sigstore) binding is implemented to strengthen them.

\section{Conclusion}
This disclosure establishes, as of 2026-06-01, public prior art for the PURIQ master
formula, the 13-axis Yuyay gate, and the Khipu DAG receipt architecture, with honest
demarcation of proved vs.\ open claims. Any later third-party attempt to patent these
disclosed methods must contend with this timestamped record.

\section*{Acknowledgment / Provenance}
All numbers herein are drawn from SZL's audited corpus (Doctrine v11 LOCKED, lutar-lean
commit \texttt{c7c0ba1}) and a real load test executed 2026-06-01. Honest caveats are
preserved per the Zero-Bandaid Law. Signed: \textbf{Yachay}.

\begin{thebibliography}{9}
\bibitem{zenodo} SZL Holdings, ``Ouroboros thesis (concept chain),'' Zenodo,
DOI: 10.5281/zenodo.19944926. \url{https://doi.org/10.5281/zenodo.19944926}
\bibitem{a11oy} SZL Holdings, ``a11oy (uds-v0.3.0),'' Zenodo,
DOI: 10.5281/zenodo.20451991.
\bibitem{lutar} SZL Holdings, ``lutar-lean v18.0.0,'' Zenodo,
DOI: 10.5281/zenodo.20434308.
\bibitem{slsa} SLSA, ``Supply-chain Levels for Software Artifacts (SLSA), Level 1.''
\url{https://slsa.dev}
\bibitem{ieee} M. Shell, ``IEEEtran document class,'' v1.8b, 2015.
\end{thebibliography}

\end{document}
